%% Section 7: Open Questions and Future Directions
%% To be \input{} into the main document.

\section{Open Questions and Future Directions}
\label{sec:open}

We organise the questions suggested by this paper into three tiers,
roughly ordered by accessibility: problems that could be attacked with
existing tools, problems that require new ideas, and problems that
call for significant new theory.


\subsection{Immediately tractable}

\begin{enumerate}[label=\textbf{Q\arabic*.}, ref=Q\arabic*, leftmargin=3em]

\item \label{q:horizontal}
\textbf{Complete the phase diagram horizontally.}
The categorical phase diagram of Section~\ref{sec:hierarchy} has
populated cells in every column of the solvable row (Schur,
Grothendieck, quantum $K$, Hall-Littlewood) and in the classical
column of the coboundary row (Demazure crystals).  The remaining
coboundary cells---K-theoretic, quantum K-theoretic, and spin models
at the quasi-solvable level---are open.  Yang's Demazure crystal
model~\cite{yang-2025} covers the classical coboundary column; the
K-theoretic and spin columns at this level remain to be constructed.
Concretely: which lattice models compute coboundary-level structure
constants for Grothendieck polynomials and Hall-Littlewood polynomials?
Do they exhibit the factored (column-level) YBE characteristic of
quasi-solvability?

\item \label{q:typeCD}
\textbf{Type C/D phase diagram.}
Everything in this paper concerns type~$A$---the general linear group
and its symmetric functions.  Azenhas, Feller, and
Torres~\cite{azenhas-feller-torres-2022} construct the symplectic
cactus group action on type~$C$ crystals, discovering additional
relations beyond those present in type~$A$.  What does the
integrability hierarchy look like for $\mathrm{Sp}_{2n}$?  The
quasi-solvable lattice models of Buciumas and
Scrimshaw~\cite{buciumas-scrimshaw} for types $B$, $C$, and $D$
provide the integrable side of the story; the categorical side---what
replaces the cactus group, whether the Hecke transition algebra
admits a type $C$ analogue---is largely unexplored.

\item \label{q:cactus-splitting}
\textbf{$J_n$ non-splitting for all~$n$.}
The cactus group $J_n$ fits into a short exact sequence
$1 \to PJ_n \to J_n \to S_n \to 1$,
where $PJ_n$ is the \emph{pure cactus group}.  For $n = 3$, we proved
by four independent methods (including the computation
$H^2(S_3, \mathbb{Z}^{\mathrm{sgn}}) \cong \mathbb{Z}/3\mathbb{Z}$)
that this sequence does not split.  Non-splitting is the algebraic
counterpart of the geometric fact that the real moduli space
$\overline{M}_{0,n+1}(\mathbb{R})$ is not a trivial $S_n$-bundle.

For $n \geq 4$, the pure cactus group is non-abelian, so the
classical group cohomology approach (which requires abelian
coefficients) does not directly apply.  The most promising route is
topological: Devyatov's presentation~\cite{devyatov-2015} of the
fundamental group of $\overline{M}_{0,n+1}(\mathbb{R})$ expresses the
pure cactus group in terms of explicit generators and relations.  A
careful analysis of these relations---perhaps combined with the
non-abelian cohomology $H^1(S_n, \mathrm{Out}(PJ_n))$---should settle
the non-splitting question for all~$n$.

\end{enumerate}


\subsection{Medium-term}

\begin{enumerate}[label=\textbf{Q\arabic*.}, ref=Q\arabic*, leftmargin=3em]
\setcounter{enumi}{3}

\item \label{q:spectral-parameter}
\textbf{$\kappa$ as spectral parameter.}
Tikhonov~\cite{tikhonov-2026} showed that the limiting permuton of a
random sorting network is governed by the invariant
$\kappa = (1 - qp)/(1-p)$, and that the M\"obius transformation
$p' = p(1-q)/(1-qp)$ relates the braided ($q > 0$) and coboundary
($q = 0$) regimes.  This transformation is an element of
$\mathrm{PGL}(2)$, acting on the probability parameter $p$ with the
Hecke parameter $q$ as coefficient.

The same group $\mathrm{PGL}(2)$ acts on the spectral parameter $u$
in the Yang-Baxter equation: the R-matrix $R(u)$ of a rational
model transforms under $u \mapsto (au + b)/(cu + d)$ by
conjugation.  Is Tikhonov's M\"obius transformation the \emph{same}
$\mathrm{PGL}(2)$ element that relates the spectral parameters of the
braided and coboundary R-matrices?  An affirmative answer would give a
quantitative, parameter-level avatar of the forgetful functor
$\mathrm{Braided} \to \mathrm{Coboundary}$ from
Section~\ref{sec:hierarchy}.

\item \label{q:cactus-yang}
\textbf{Cactus action on Yang's boundary conditions.}
Yang's Demazure crystal model~\cite{yang-2025} is parameterised by a
permutation $w \in S_r$, which enters as the boundary condition on the
right edge of the lattice.  The cactus group $J_r$ acts on Demazure
crystals via compositions of Sch\"utzenberger
involutions~\cite{henriques-kamnitzer-2006}.  Does the cactus action
on the crystal correspond to a permutation of Yang's boundary
conditions---that is, does the cactus generator $s_{[i,j]}$ send the
boundary condition $w$ to $w \cdot w_0^{[i,j]}$, where $w_0^{[i,j]}$
is the longest element of $S_{\{i, \ldots, j\}}$?  If so, the
coboundary commutor would have an explicit lattice-model
manifestation: it permutes the boundary data, and the resulting
bijection on configurations is the lattice-model avatar of the
Sch\"utzenberger involution.

\item \label{q:specialization-moduli}
\textbf{Specialization moduli.}
Fan, Guo, and Xiong~\cite{fan-guo-xiong} proved that KTW puzzles and
bumpless pipe dreams are not independent models but rather
\emph{specialisations} of a single lattice model: different choices of
boundary conditions and spectral parameters in one unified vertex model
recover both combinatorial objects.  What is the \emph{moduli space} of
all such specialisations?  Is it related to the space of factored
integrability conditions (the locus where $R_{\mathrm{row}}$ and
$R_{\mathrm{col}}$ independently satisfy the YBE)?  A complete
description of this moduli space would answer, in a precise sense,
the question ``how many essentially different combinatorial models for
the LR coefficients exist?''

\end{enumerate}


\subsection{Long-term}

\begin{enumerate}[label=\textbf{Q\arabic*.}, ref=Q\arabic*, leftmargin=3em]
\setcounter{enumi}{6}

\item \label{q:shuffle-full}
\textbf{Shuffle algebra for the full phase diagram.}
The Garbali-Negut isomorphism~\cite{garbali-negut-2025} connects
shuffle algebras to Yangians at the rational level, which governs the
classical (cohomological) column of the phase diagram.  Does this
isomorphism extend to the trigonometric level (K-theoretic column,
where the relevant quantum group is the quantum affine algebra
$U_q(\widehat{\mathfrak{sl}}_n)$) and the elliptic level (where the
quantum group is Felder's elliptic quantum
group~\cite{felder-1995})?  The Shibukawa-Ueno
hierarchy~\cite{shibukawa-ueno}---elliptic, trigonometric,
rational---classifies R-matrices by the function field in which they
live, and is \emph{orthogonal} to the braided/coboundary/symmetric
hierarchy of Section~\ref{sec:hierarchy}.  The intersection of these
two hierarchies would give a three-dimensional phase diagram: one axis
for deformation type, one for integrability level, and one for the
analytic class of the R-matrix.

\item \label{q:tqft}
\textbf{TQFT for puzzle bundles.}
The failure of the Zinn-Justin tetrahedron equation at the
quasi-solvable level (Section~\ref{sec:groupoid-hierarchy}) defines a
curvature 2-form on the ``puzzle bundle'' over the higher Bruhat order
$B(n,2)$.  This curvature measures the obstruction to promoting
column-level integrability to full vertex-level integrability, and its
holonomy around closed loops encodes the non-abelian structure of the
pure cactus group.  Is there a topological quantum field theory whose
partition function computes LR coefficients, with the curvature
encoding the quasi-solvable level?  Such a TQFT would provide a
geometric home for the categorical phase diagram: the braided level
would correspond to a flat connection (zero curvature, trivial
holonomy), the coboundary level to a curved connection with discrete
holonomy, and the symmetric level to the trivial bundle.

\item \label{q:kpz}
\textbf{KPZ universality for Schubert calculus.}
Morales, Panova, Petrov, and
Yeliussizov~\cite{morales-panova-petrov-yeliussizov} connect pipe
dreams to the totally asymmetric simple exclusion process (TASEP) and
Tracy-Widom distributions, placing classical Schubert calculus in the
KPZ universality class.  Does this connection extend to K-theoretic or
quantum K-theoretic Schubert calculus?  The K-theoretic structure
constants are \emph{signed}, which complicates the probabilistic
interpretation, but Tikhonov's permuton results~\cite{tikhonov-2026}
suggest that the probabilistic side is tractable whenever the Hecke
parameter $q$ is real.  A KPZ universality theorem for Grothendieck
polynomials would link the algebraic geometry of flag varieties to the
statistical mechanics of interface growth---two subjects that have
developed largely in parallel.

\end{enumerate}

\bigskip

\noindent
\textbf{Coda.}\enspace
The central message of this paper is that the multiplicity of
combinatorial models for LR coefficients is not a historical accident
but a structural inevitability.  The models proliferate because they
are representations of a rich algebraic object---the shuffle algebra
and its quotients---organised by the integrability hierarchy into
levels of increasing symmetry.  Each level forgets something precise
(a spectral parameter, a multiplicity bundle), and each forgetting
produces a simpler but still faithful encoding of the same
nonnegative integers.  Understanding this structure does more than
explain existing combinatorics: it predicts where new models should
live (the open cells of the phase diagram), identifies the invariants
that distinguish them (the Hecke parameter $q$, the shuffle
kernel~$\zeta$), and points toward deeper explanations (geometric
Satake, TQFT, KPZ universality) that may one day subsume the entire
picture.

The LR coefficients, those modest nonnegative integers governing the
multiplication of Schur functions, have turned out to be a window onto
a remarkably deep landscape.  We have tried to draw a map of what we
can see.  The territory beyond the map---the open cells, the
conjectural connections, the questions we do not yet know how to
ask---is, we suspect, richer still.
